\chapter{Applications}
\label{ch:apps}

The design choices of \NS---a million-token context, high throughput on long
generations, and post-training focused on multi-step tool use---point to a
particular class of applications. The Nemotron~3 white paper positions Super
for collaborative agents and high-volume workloads such as IT ticket
automation~\cite{nvidia2025nemotron3}. The main application areas are outlined
below.

\section{Coding Agents}

A coding agent must read a repository, locate the cause of an issue, edit
several files and run the tests, often over dozens of steps. Each step adds
tool output to the context, so the context grows quickly and most of the work
is \emph{prefill}: processing a long input before a short action. This is the
workload for which the NVFP4 checkpoint was designed. \NS\ was trained for it
directly, with SFT data built from real GitHub issues and a dedicated
software-engineering RL stage, and it scores 60.5 on SWE-Bench with the
OpenHands harness (Table~\ref{tab:main}).

\section{Terminal and Multi-Step Task Automation}

Closely related are assistants that operate a command line: installing
software, inspecting logs, and chaining commands to complete an administrative
task. The model was trained on terminal-use data and is evaluated on Terminal
Bench. Long sessions benefit from the constant-size Mamba state, because the
cost of each new token does not climb as the session history accumulates in
the way it does for a model built only from attention layers.

\section{Tool-Using Agents and Customer Service}

Many business applications consist of a model calling external tools: querying
a database, searching the web, or invoking an API to check an order or change
a booking. \NS\ was trained with conversational and general-purpose
tool-calling data and with search and SQL datasets, and is evaluated on
TauBench (airline, retail and telecom customer-service scenarios), BrowseComp
(web search) and BIRD (text-to-SQL). High throughput matters here because such
systems serve many users at once and often run several model calls per
request.

\section{Long-Document Analysis}

With a context of one million tokens, entire contracts, sets of research
papers, financial reports or codebases can be given to the model at once,
instead of being cut into fragments and retrieved piecemeal. The RULER score
of 91.6 at one million tokens indicates that the model can actually locate and
use information at that range, rather than merely accept it as input.

\section{Reasoning with Controllable Effort}

The model supports three reasoning modes---off, regular and low-effort---which
can be combined with a reasoning budget set at inference time. An application
can thus spend long reasoning traces on hard problems and answer simple
requests quickly, trading accuracy against cost request by request.

\section{On-Premises and Customised Deployment}

Because the weights, the quantized checkpoints and much of the training data
are openly released, an organisation can run the model on its own hardware and
fine-tune it on its own data. This is important where data cannot leave the
organisation, for example in finance, healthcare or government. The FP8 and
NVFP4 checkpoints target the Hopper and Blackwell GPU generations
respectively, and the built-in MTP heads provide speculative decoding without
a separate draft model.
